\title{Direct Preference Optimization: Your Language Model is Secretly a Reward Model}

\begin{abstract}
Although large unsupervised language models (LMs) acquire extensive world knowledge and develop some reasoning abilities, fine-grained management of their outputs remains challenging due to their fully unsupervised training regimes. Approaches to enhance this controllability typically involve gathering human judgments regarding the relative merit of generated responses and subsequently fine-tuning the base LM to conform to these judgments—frequently through reinforcement learning from human feedback (RLHF). RLHF, however, is often intricate and can be unstable: it consists of first training a reward model to represent human preferences and then applying reinforcement learning to adapt the LM to maximize this learned reward while maintaining proximity to the initial model parameters. In this work, we propose a novel way to parameterize the reward model within RLHF that permits direct, closed-form computation of the corresponding optimal policy, reducing the RLHF objective to a standard classification loss. We term this approach \textit{Direct Preference Optimization} (DPO). DPO is computationally efficient, reliable, and obviates the need for generating samples from the LM during fine-tuning or for extensive hyperparameter optimization. Empirically, we find that DPO is capable of fine-tuning LMs to match or outperform existing techniques in adhering to human preferences. Specifically, DPO achieves superior sentiment control in language generation compared to RLHF with PPO, and is competitive or better with respect to summarization and single-turn dialogue tasks—all while offering a substantially simpler training and implementation process.
\end{abstract}

\section{Introduction}
Large-scale unsupervised language models (LMs), when exposed to extensive datasets, often demonstrate unexpected capabilities~\citep{chowdhery2022palm, brown2020language, touvron2023llama, bubeck2023sparks}. Despite these achievements, the underlying data is produced by humans with diverse intentions, competencies, and objectives. This diversity can result in models learning behaviors that are not always desirable; for instance, although it is beneficial for an AI programming assistant to \textit{recognize} frequent coding errors so as to rectify them, we would ultimately prefer the model to favor the (possibly uncommon) examples of expert coding embedded in its training corpus when generating code. Likewise, while it is advantageous for a language model to be \textit{informed} of a widespread misconception (shared by, say, half the population), we certainly do not want it to propagate this erroneous belief in 50\% of relevant queries. Thus, carefully shaping the model’s \emph{preferred outputs and conduct} from its extensive \textit{repository of knowledge and capabilities} is essential for constructing AI that is robust, high-performing, and governable \citep{ouyang2022training}. Although reinforcement learning (RL) has been the standard for aligning LMs with human preferences, we will demonstrate that the RL-based objectives employed by current approaches are in fact equivalent to a simple binary cross-entropy loss, offering a more straightforward alternative to preference optimization.

Broadly, present-day approaches impart desired traits to language models by leveraging specifically designed collections of human preference data that encapsulate behaviors regarded as helpful and safe by people. This phase of preference alignment typically follows large-scale, unsupervised pre-training on diverse text corpora. While one basic strategy is to fine-tune using supervised learning on human-crafted examples of desirable responses, reinforcement learning from human or AI feedback (RLHF/RLAIF; \citep{christiano2017deep,bai2022constitutional}) has proven particularly effective. RLHF methodologies learn a reward function from annotated human preferences and subsequently use RL algorithms to tune the model so that its outputs attain higher rewards, all while constraining deviation from the pre-trained baseline. Despite yielding models with state-of-the-art abilities in conversation and programming, the RLHF workflow is notably more involved than supervised fine-tuning; it requires training several models and repeated policy sampling during the training process, which results in substantial computational demands.

This study presents a method for optimizing language models in accordance with human preferences, foregoing the need for explicit reward models or reinforcement learning techniques. We introduce \textit{Direct Preference Optimization} (DPO), a novel algorithm that targets the same underlying objective as conventional RLHF methods—maximizing reward subject to a KL-divergence penalty—yet offers a more accessible implementation and training procedure. Conceptually, DPO adjusts the model by amplifying the relative log likelihood of preferred responses compared to less preferred ones, incorporating a flexible, example-specific importance weighting mechanism. This approach mitigates the model collapse often observed with simpler probability-ratio objectives. DPO, like other methods, makes use of a theoretical preference framework (for example, the Bradley-Terry model; \cite{bradley1952rankanalysis}) to quantify how well a reward signal matches observed preference judgments. The key distinction lies in the fact that traditional pipelines leverage this framework to establish a preference loss for training a reward model and then update a policy to maximize the learned reward, whereas DPO reformulates the problem, expressing the preference loss directly in terms of the policy via a variable transformation. When equipped with annotated data capturing human preferences among model outputs, DPO efficiently optimizes the policy using a standard binary cross-entropy objective, ultimately yielding the best-fit policy relative to an implicit reward informed by the preference data.

Our principal contribution is Direct Preference Optimization (DPO), which provides a straightforward alternative to RL-based methods for preference-driven language model training. Through empirical evaluation, we demonstrate that DPO performs on par with or surpasses established approaches—including PPO-based RLHF—across various preference-driven tasks, such as sentiment control, summarization, and conversational response, even when scaling to language models of up to 6B parameters.

\section{Related Work}

As self-supervised language models increase in scale, they are capable of handling certain tasks in a zero-shot manner \citep{radford2019language}, or by using few-shot prompting strategies \citep{gpt3,megatron,chowdhery2022palm}. Nonetheless, refining these models by fine-tuning on curated datasets comprised of explicit instructions and human-authored completions can lead to marked gains in task-specific performance and better adherence to user intentions \citep{mishra-etal-2022-cross,sanh2022multitask,chung2022scaling,thoppilan2022lamda}. This technique, known as `instruction-tuning,' not only boosts the general utility of large language models (LLMs) but also empowers them to successfully process instructions not present during the tuning phase \citep{chung2022scaling}. In practice, although collecting precise human demonstrations for instruction tuning is resource-intensive, gathering \textit{relative} human assessments of output quality tends to be more feasible. As a result, more recent approaches have focused on fine-tuning LLMs with datasets capturing human preferences, leading to notable progress in applications such as translation \citep{kreutzer-etal-2018-reliability}, summarization \citep{stiennon2022learning,ziegler2020finetuning}, narrative generation \citep{ziegler2020finetuning}, and instruction-following tasks \citep{ouyang2022training,ramamurthy2023is}. These systems typically begin by training a reward model, often within the Bradley-Terry framework \citep{bradley1952rankanalysis}, to fit the collected preference data. Subsequently, the LLM is further optimized through reinforcement learning, using techniques like REINFORCE \citep{williams1992reinforce}, proximal policy optimization (PPO; \cite{schulman2017proximal}), or related methods \citep{ramamurthy2023is}, with the aim of maximizing the estimated rewards. Related research streams use instruction-tuned LLMs—already augmented with human feedback—to generate synthetic preference data emphasizing particular dimensions such as safety or non-harmfulness \citep{bai2022constitutional}, relying solely on minimally structured guidance (e.g., text-based rubrics) for LLM-generated evaluations. These advances exemplify the synthesis of research on reinforcement learning-based language model training for diverse end-goals~\citep{Ranzato2015SequenceLT,paulus2018a,wu2018learning} and studies investigating algorithms for learning from human preferences \citep{christiano2017deep,kupcsik2018learning}. Notwithstanding the advantages of leveraging relative human preferences, scaling reinforcement learning approaches for LLM fine-tuning remains a complex practical hurdle; this paper proposes a methodologically sound alternative to optimize for relative preferences without resorting to reinforcement learning.

Beyond language-specific applications, the study of learning policies from preference feedback has been explored within both bandit and reinforcement learning frameworks, leading to the development of multiple methodologies. In contextual bandit problems, when learning is based on preferences or rankings of actions instead of explicit reward signals, the setting is termed as contextual dueling bandits (CDB; \cite{yue2012karmed,dudik2015contextual}). Without direct access to absolute rewards, theoretical treatments of CDBs introduce the concept of the \textit{von Neumann winner}—a policy whose expected rate of outperforming any alternative policy meets or exceeds 50\% \citep{dudik2015contextual}. Unlike CDBs, where preference annotations are collected in an online manner, learning from human preferences most often utilizes a pre-existing, static dataset of action pairs labeled with preferences \citep{yan2022human}. In parallel, \textit{preference-based RL} (PbRL) algorithms obtain feedback in the form of binary preferences, which are presumed to arise from an unspecified underlying ‘scoring’ or reward function \citep{BusaFekete2014,ruiz2023dueling}. A variety of PbRL methods have been introduced, some of which allow leveraging off-policy preference data. Typically, these approaches involve an initial step of estimating this latent reward or scoring function (i.e., constructing a reward model), which is then used for policy optimization \citep{jain2013learning,BusaFekete2014,christiano2017deep,sadigh2017active,kupcsik2018learning}. In contrast, our method proposes a unified, one-stage framework that directly trains a policy to align with the observed preference data.

\section{Preliminaries}\label{section:prelims}

We provide an overview of the RLHF process outlined by \citeauthor{ziegler2020finetuning}, as well as subsequent works \citep{stiennon2022learning, bai2022training, ouyang2022training}. This approach conventionally proceeds through three primary stages: (1) supervised fine-tuning (SFT), (2) the collection of preference data along with reward model training, and (3) reinforcement learning-based optimization.

\textbf{SFT}: The RLHF workflow generally starts with adapting a large, pre-trained language model through supervised learning on task-specific, high-quality data, producing the SFT model $\pi^\text{SFT}$. This phase targets downstream applications such as dialogue or summarization.

\textbf{Reward Modelling Phase}: During the second stage, prompts $x$ are submitted to the SFT model, resulting in the generation of pairs of candidate responses $(y_1, y_2)\sim \pi^\text{SFT}(y \mid x)$. Human annotators then assess these pairs, expressing a preference, denoted as $y_w\succ y_l \mid x$, with $y_w$ and $y_l$ indicating the chosen and non-chosen responses, respectively, for each prompt. Although these preferences are governed by an underlying reward mechanism $r^*(y, x)$ that remains hidden, several modeling strategies can be applied. The Bradley-Terry (BT) model \cite{bradley1952rankanalysis} is frequently employed for preference modeling, though the more general Plackett-Luce ranking model \citep{plackett1975analysis, luce2012individual} can be used when data contains multi-way rankings. Under the BT formulation, the human-induced preference distribution $p^*$ can be described as follows:

Given a static dataset of preference comparisons $\mathcal{D}=\bigl\{x^{(i)}, y_w^{(i)}, y_l^{(i)}\bigr\}_{i=1}^N$ sampled from $p^*$, one can parameterize a reward model $r_{\phi}(x, y)$ and fit its parameters by maximizing likelihood. When reinterpreting the problem as a binary classification task, the negative log-likelihood objective is given by:

\begin{equation}\label{eq:reward_model}
    \mathcal{L}_R(r_{\phi}, \mathcal{D}) = -\mathbb{E}_{(x, y_w, y_l)\sim \mathcal{D}}\bigl[\log \sigma(r_{\phi}(x, y_w)- r_{\phi}(x, y_l))\bigr]
\end{equation}

where $\sigma$ denotes the logistic sigmoid function. In typical language modeling applications, $r_{\phi}(x, y)$ is initialized from the supervised fine-tuning (SFT) model $\pi^\text{SFT}(y \mid x)$ and augmented by appending a linear readout on the final transformer hidden states to yield a scalar reward output~\cite{ziegler2020finetuning}. To control reward variance, earlier works have normalized reward values so that $\mathbb{E}_{x,y\sim \mathcal{D}}\left[r_\phi(x, y)\right] = 0$ holds across all $x$.

\textbf{RL Fine-Tuning Phase}: In the reinforcement learning stage, the trained reward model is employed to supply feedback for optimizing the language model. Following established work~\citep{jaques2017sequence, jaques2020human}, the learning objective takes the form:

\begin{equation}\label{eq:RL}
\max_{\pi_{\theta}}  \mathbb{E}_{x\sim \mathcal{D}, y\sim \pi_{\theta}(y \mid x)}\bigl[r_{\phi}(x, y)\bigr] - \beta\mathbb{D}_{\textrm{KL}}\bigl[\pi_{\theta}(y\mid x)\mid \mid \pi_\text{ref}(y\mid x)\bigr],
\end{equation}

Here, $\beta$ modulates how much the policy $\pi_\theta$ is permitted to stray from the reference policy $\pi_\text{ref}$, typically the initial supervised fine-tuned policy $\pi^\text{SFT}$. In standard implementations, $\pi_\theta$ is also initialized from $\pi^\text{SFT}$. This KL regularization is critical—it constrains the model within a region where the reward predictor is reliable, while preserving diversity in generation and guarding against mode collapse. As the generation process is discrete, direct gradient-based optimization is not feasible, and reinforcement learning methods are used. Customarily, the reward for policy learning is composed as $r(x, y) = r_{\phi}(x, y) -\beta (\log \pi_{\theta}(y\mid x) - \log \pi_\text{ref}(y\mid x))$, which is then maximized with PPO~\cite{schulman2017proximal} according to previous work~\cite{ziegler2020finetuning, stiennon2022learning, bai2022training, ouyang2022training}.

\section{Direct Preference Optimization}\label{sec:DPO}

Recognizing the practical difficulties of applying reinforcement learning to large-scale language model fine-tuning, we seek a more direct alternative for optimizing policies from human preferences. Whereas classical RLHF methods construct an explicit reward function for downstream RL training, we instead identify a parameterization of the reward model that exposes its corresponding optimal policy in closed form—obviating the need for any RL loop. Our central idea is to apply an analytic relationship between rewards and their optimal policies, allowing us to recast the reward-learning loss in terms of policies themselves. This change of variables eliminates the requirement to independently fit a reward model, yet maintains alignment with human preference structures such as the Bradley-Terry framework. Effectively, the policy model fulfills the dual roles of both the language model and the (implicit) reward representation.

\textbf{DPO Objective Derivation.} Our analysis commences from the standard RL objective as utilized in previous studies, specifically Eq.~\ref{eq:RL}, using an arbitrary reward function $r$. As discussed in earlier literature~\citep{peters2007reinforcement, peng2019advantage, korbak2022reinforcement, go2023aligning}, it is direct to verify that the policy which maximizes the KL-regularized reward in Eq.~\ref{eq:RL} admits the following structure:

\begin{equation}\label{eq:op_policy}
    \pi_r(y\mid x) = \frac{1}{Z(x)}\pi_\text{ref}(y\mid x)\exp\left(\frac{1}{\beta}r(x, y)\right),
\end{equation}

with $Z(x) =\sum_{y}\pi_\text{ref}(y\mid x)\exp\left(\frac{1}{\beta}r(x, y)\right)$ serving as the normalization constant (partition function). A detailed walkthrough can be found in Appendix \ref{app:derivation1}. In practice, however, even when $r_{\phi}$ serves as the MLE approximation of the true reward $r^*$, calculating the partition function $Z(x)$ remains computationally challenging~\citep{korbak2022reinforcement, go2023aligning}. This difficulty limits the direct applicability of the formulation. Nonetheless, by manipulating Eq.~\ref{eq:op_policy}, the reward function can be rewritten solely in terms of the optimal policy $\pi_r$, the baseline policy $\pi_\text{ref}$, and the normalization factor $Z(\cdot)$. Concretely, after taking logarithms of both sides of Eq.~\ref{eq:op_policy} and rearranging, we have:

\begin{equation}\label{eq:main_eq}
    r(x,y) =\beta \log \frac{\pi_r(y\mid x)}{\pi_\text{ref}(y\mid x)} + \beta \log Z(x).
\end{equation}

Applying this change of variables for both the authentic reward $r^*$ and its optimal policy $\pi^*$, we observe that the Bradley-Terry framework only leverages the reward differences between pairs of completions: ${p^*(y_1 \succ y_2 \mid x) = \sigma(r^*(x, y_1) - r^*(x, y_2))}$. Plugging the representation from Eq.~\ref{eq:main_eq} into the preference function (Eq.~\ref{eq:bradley-terry}), the normalization $Z(x)$ cancels out, leaving the probability of user preference defined only through $\pi^*$ and $\pi_\text{ref}$. Therefore, within the Bradley-Terry modeling, the RLHF-optimal policy $\pi^*$ fulfills:

\begin{equation}\label{eq:objective}
    p^*(y_1\succ y_2 \mid x)=\frac{1}{1 + \exp\left(\beta \log \frac{\pi^*(y_2\mid x)}{\pi_\text{ref}(y_2\mid x)} - \beta \log \frac{\pi^*(y_1\mid x)}{\pi_\text{ref}(y_1\mid x)}\right)}
\end{equation}

A comprehensive proof can be found in Appendix~\ref{app:derivation2}. Although this derivation relies on the Bradley-Terry model, the argument extends straightforwardly to broader ranking models such as the Plackett-Luce family~\citep{plackett1975analysis, luce2012individual}, elaborated upon in Appendix~\ref{app:plackett_luce_models}.

Since we have an explicit likelihood over human preference annotations as a function of the policy parameters, we can directly optimize a parameterized policy $\pi_\theta$ using maximum likelihood. Following the same principle as in the reward modeling paradigm (see Eq.~\ref{eq:reward_model}), the resulting policy optimization criterion can be written as:

\begin{equation}\label{eq:optimum_model}
    \mathcal{L}_\text{DPO}(\pi_{\theta}; \pi_\text{ref}) = -\mathbb{E}_{(x, y_w, y_l)\sim \mathcal{D}}\left[\log \sigma \left(\beta \log \frac{\pi_{\theta}(y_w\mid x)}{\pi_\text{ref}(y_w\mid x)} - \beta \log \frac{\pi_{\theta}(y_l\mid

In this formulation, we estimate an implicit reward through a different parameterization such that the optimal policy corresponds directly to $\pi_\theta$. Furthermore, our approach mirrors fitting a reparameterized Bradley-Terry model, granting the method certain theoretical guarantees—including consistency under appropriate assumptions regarding the distribution of preference data \cite{bong2022generalized}. A more detailed examination of the theoretical foundations of DPO and its connections to related work is provided in Section~\ref{sec:theory}.

\textbf{Mechanistic insight into the DPO update:} To gain intuition into the inner workings of DPO, we can inspect the gradient of its loss function, $\mathcal{L}_\text{DPO}$. The gradient with respect to model parameters $\theta$ is given by:

\begin{multline*}\label{eq:gradient}
    \nabla_\theta \mathcal{L}_\text{DPO}(\pi_\theta;\pi_\text{ref}) = \\ -\beta\mathbb{E}_{(x, y_w, y_l) \sim \mathcal{D}} \bigg[\underbrace{\sigma(\hat{r}_\theta(x, y_l) - \hat{r}_\theta (x, y_w))}_\text{assigns greater weight when the reward estimate misorders the preference}\bigg[\underbrace{\nabla_\theta\log \pi(y_w \mid x)}_\text{increase probability of $y_w$} - \underbrace{\nabla_\theta\log\pi(y_l \mid x)}_\text{reduce probability of $y_l$}\bigg]\bigg],
\end{multline*}

where the reward proxy is defined as $\hat{r}_\theta(x, y) = \beta \log \frac{\pi_\theta(y \mid x)}{\pi_\text{ref}(y \mid x)}$, drawing on both the current language model $\pi_\theta$ and a reference model $\pi_\text{ref}$ (elaborated in Section~\ref{sec:theory}). The practical effect of this gradient is to push the model to favor the more desirable completions $y_w$ while diminishing the likelihood assigned to the less favored completions $y_l$. Crucially, the degree to which an example influences the update is determined by how strongly the implicit reward model $\hat{r}_\theta$ incorrectly prefers the disfavored option, modulated by $\beta$, thus reflecting the severity of the ordering mistake as well as the strictness of the KL regularization. Empirical results highlight the necessity of this weighting: omitting the weighting term leads to inferior outcomes, with models at risk of collapse as evidenced in Appendix Table~\ref{tab:unlikelihood_generations}.

\textbf{DPO overview.} The typical workflow for DPO consists of the following steps: 1) For each prompt $x$, generate candidate completions $y_1, y_2 \sim \pi_\text{ref}(\cdot \mid x)$, collect human preference annotations to build an offline preference dataset $\mathcal{D} = \{x^{(i)}, y_w^{(i)}, y_l^{(i)}\}_{i=1}^N$; 2) update the language model $\pi_\theta$ by minimizing the objective $\mathcal{L}_\text{DPO}$, conditioned on $\pi_\text{ref}$, $\mathcal{D}$, and a chosen $\beta$.  
In practice, preference datasets that have already been released are often preferred to avoid the costly process of new sample generation and preference elicitation. If these datasets were produced using $\pi^\text{SFT}$, we set the reference policy $\pi_\text{ref}$ equal to $\pi^\text{SFT}$ whenever possible. When such a policy is not accessible, we instead fit $\pi_\text{ref}$ by maximizing the likelihood over the set of preferred completions ${(x, y_w)}$, i.e., setting ${\pi_\text{ref} = \argmax_{\pi}\mathbb{E}_{x, y_w \sim \mathcal{D}}\left[\log \pi(y_w \mid x)\right]}$. This approach helps reduce the mismatch between the (unknown) true reference distribution and the proxy $\pi_\text{ref}$ utilized in DPO. Details on the practical implementation and chosen hyperparameters are further described in Appendix~\ref{app:implementation}.

\section{Theoretical Analysis of DPO} In this section, we further clarify the intuition behind DPO, present supporting theoretical results, and discuss how the strengths of DPO compare to actor-critic methods popular in RLHF, such as PPO~\cite{schulman2017proximal}, addressing known limitations in those approaches.

\label{sec:theory}

\subsection{Your Language Model Is Secretly a Reward Model} DPO removes the need for explicit reward modeling and separate reinforcement learning by reframing policy optimization as a single maximum likelihood task. Importantly, the loss function in Eq. \ref{eq:main_eq} mirrors the structure of the Bradley-Terry model when parametrized with a reward function $r^*(x, y) = \beta \log\frac{\pi^*_\theta(y \mid x)}{\pi_\text{ref}(y \mid x)}$. Optimization is then carried out with respect to the parameterized model $\pi_{\theta}$, equivalent to optimizing a reward model as in Eq. \ref{eq:reward_model} but under a variable transformation. Here, we develop the theoretical foundations for this reparameterization, demonstrate that it does not restrict the set of achievable reward models, and prove that it permits recovery of the optimal policy. We start by introducing an equivalence relationship between reward functions.

\begin{definition}
Two reward functions $r(x, y)$ and $r'(x, y)$ are called equivalent if and only if ${r(x, y)-r'(x, y) = f(x)}$ holds for some function $f$.
\end{definition}

This relation is easily verified to be an equivalence relation, thereby partitioning the collection of reward functions into distinct equivalence classes. We now present two supporting lemmas:

\begin{lemma}\label{lemma:same_prefrence} Within the Plackett-Luce preference model—including the Bradley-Terry formulation—any pair of reward functions from the same equivalence class generates identical preference distributions.
\end{lemma}

\begin{lemma}\label{lemma:same_policy}
    Any two reward functions from a common equivalence class yield the same optimal policy for the constrained RL setting.
\end{lemma}

The proofs follow directly and are provided in Appendix \ref{app:lemma1}. The first lemma reiterates a classic under-specification property found in Plackett-Luce models \cite{plackett1975analysis}. This lack of identifiability typically necessitates the enforcement of extra constraints to ensure meaningful MLE estimation from Eq. \ref{eq:reward_model} \cite{bong2022generalized}. According to the second lemma, all reward functions within the same equivalence class produce an identical optimal policy. Therefore, for our overall goal, it suffices to identify any reward function within the optimal equivalence class. In Appendix~\ref{app:thm1}, we formally establish the following theorem:

\begin{theorem}\label{thm:main}
    Assuming mild technical conditions, every reward equivalence class compatible with Plackett-Luce models (including Bradley-Terry) can be parameterized by ${r(x, y) = \beta \log \frac{\pi(y\mid x)}{\pi_\text{ref}(y\mid x)}}$, for some $\pi(y\mid x)$ and a reference distribution $\pi_\text{ref}(y \mid x)$.
\end{theorem}

\begin{sproof}
    Take any reward function $r(x, y)$ which, via Eq. \ref{eq:op_policy}, defines its optimal policy $\pi_r(y \mid x)$. We will demonstrate that a member of the equivalence class of $r$ can be written in the reparameterized form above. Define the mapping $f$ as
\begin{equation}
    f(r; \pi_\text{ref}, \beta)(x, y) = r(x, y) - \beta\log\sum_{y}\pi_\text{ref}(y\mid x)\exp\left(\frac{1}{\beta}r(x, y)\right)
\end{equation}
The map $f$ effectively normalizes the reward by subtracting the log partition function of $\pi_r$, which depends only on $x$. Consequently, $f(r; \pi_\text{ref}, \beta)(x, y)$ remains in the equivalence class of $r(x, y)$. Moreover, if we substitute $r$ by the right-hand side of Eq.~\ref{eq:main_eq} (applicable for any $r$), it follows that $f(r; \pi_\text{ref}, \beta)(x, y) = \beta \log \frac{\pi_r(y\mid x)}{\pi_\text{ref}(y\mid x)}$. Therefore, this mapping $f$ selects a representative of the equivalence class of $r$ in the desired parameterization, showing that our reparameterization retains complete generality.
\end{sproof}

Theorem~\ref{thm:main} can also be interpreted as identifying, within each equivalence class, the specific reward function that the DPO reparameterization selects—namely, the function that fulfills:

\begin{equation}\label{eq:lag_p}
     \sum_{y}\underbrace{\pi_\text{ref}(y\mid x)\exp\left(\frac{1}{\beta}r(x, y)\right)}_{=\pi(y\mid x)\text{, using Thm.~\ref{thm:main} reparam.}} = 1,
\end{equation}

In other words, $\pi(y\mid x)$ constitutes a legitimate probability distribution, as all probabilities are non-negative and sum to one. Examining Eq.~\ref{eq:op_policy}, it becomes evident that Eq.~\ref{eq:lag_p} corresponds to the partition function associated with the optimal policy derived from the reward function $r(x, y)$. The central insight underlying the DPO algorithm is the ability to impose targeted constraints on the Plackett-Luce preference models—which include the Bradley-Terry model—so that the expressive power of the reward models remains unaffected, while also ensuring that the optimal policy described in Eq. \ref{eq:op_policy} remains analytically tractable across all prompts $x$.

\subsection{Instability of Actor-Critic Algorithms} Our framework further allows us to analyze why standard actor-critic algorithms applied in RLHF, such as PPO, may exhibit instability. Adhering to the RLHF pipeline, we concentrate on the RL fine-tuning phase delineated in Section \ref{section:prelims}. Connections can be made to the control as inference paradigm~\cite{levine2018reinforcement}, particularly as it pertains to the constrained RL setting introduced in \ref{eq:RL}. With a parameterized policy model $\pi_{\theta}(y\mid x)$, the objective is to minimize $\mathbb{D}_{\text{KL}}[\pi_{\theta}(y|x) \mid \mid \pi^*(y\mid x)]$, where $\pi^*$ represents the optimal policy in Eq. \ref{eq:optimum_model} as determined by the reward function $r_{\phi}(y, x)$. Basic algebraic manipulation yields the following optimization criterion:

\begin{equation}\label{eq:AC}
    \max_{\pi_{\theta}}\mathbb{E}_{\pi_{\theta}(y\mid x)}\bigg[\underbrace{r_{\phi}(x, y) -\beta\log\sum_{y}\pi_\text{ref}(y\mid x)\exp\left(\frac{1}{\beta}r_{\phi}(x, y)\right)}_{f(r_{\phi}, \pi_\text{ref}, \beta)} - \underbrace{\beta\log\frac{\pi_{\theta}(y\mid x)}{\pi_\text{ref}(y\mid x)}}_{\text{KL}}\bigg]
\end{equation}

This objective is identical to that targeted in earlier studies 
\citep{ziegler2020finetuning, stiennon2022learning, bai2022training, ouyang2022training}, where a DPO-equivalent reward is applied to the reward function class $r_{\phi}$. Within this context, the normalization component in $f(r_{\phi}, \pi_\text{ref}, \beta)$ can be seen as representing the soft value function for the reference policy $\pi_\text{ref}$. Although this normalization does not alter the final optimal solution, its absence can result in the policy gradient exhibiting substantial variance, thus destabilizing learning. One way to incorporate this normalization is through a value function learned during training, but this approach presents its own optimization challenges. Alternatively, previous approaches have adopted normalization of rewards against a baseline consisting of human completions, which effectively serves as a single-sample Monte Carlo approximation of the normalization factor. In contrast, DPO’s reparameterization framework defines a reward function that eliminates the need for such baselines. 

\section{Experiments}
Here, we empirically assess how effectively DPO trains policies from preference data alone. We start by considering a carefully controlled text-generation task, examining how DPO balances maximizing expected reward with minimizing KL-divergence to the reference policy in comparison to standard preference learning approaches like PPO. Subsequently, we investigate DPO’s performance when scaled to larger models and more complex RLHF scenarios, such as summarization and dialogue. Our findings indicate that DPO typically matches or exceeds the performance of robust baselines—including RLHF with PPO and selection of the highest-rewarded trajectory among $N$ samples—using minimal hyperparameter adjustment. We detail our experimental procedures below; further specifics are provided in Appendix~\ref{app:exp_details}.

\textbf{Tasks.} Our empirical study investigates three distinct tasks involving open-ended text generation. In every case, the models learn a policy from a preference dataset $\mathcal{D}=\bigl\{x^{(i)}, y_w^{(i)}, y_l^{(i)}\bigr\}_{i=1}^N$. For \textbf{controlled sentiment generation}, the input $x$ consists of a review prompt taken from the IMDb dataset \cite{maas-EtAl:2011:ACL-HLT2011}, with the goal that the model produces an output $y$ conveying positive sentiment. To ensure a controlled assessment in this scenario, we automatically construct preference pairs of outputs using a pre-existing sentiment classifier, selecting such that $p(\text{positive}\mid x,y_w)>p(\text{positive}\mid x,y_l)$. For SFT, GPT-2-large is fine-tuned until it converges, utilizing the training split of the IMDB dataset; additional details are presented in App~\ref{app:sentiment_details}. The \textbf{summarization} task uses $x$ as a Reddit forum post, and requires generating a summary $y$ that concisely conveys its primary content. We build upon prior approaches, utilizing the Reddit TL;DR summarization dataset \citep{volske-etal-2017-tl} alongside the human preference data collected by \citeauthor{stiennon2022learning}. Our SFT model is adapted to this task through fine-tuning on summaries composed by humans\footnote{\url{https://huggingface.co/CarperAI/openai_summarize_tldr_sft}}, employing the TRLX framework \citep{leandro_von_werra_2023_7790115} for RLHF training. Notably, the human preferences in the dataset were elicited using samples generated from a comparable, yet separately fine-tuned, SFT model as described in \citeauthor{stiennon2022learning}. For the \textbf{single-turn dialogue} task, $x$ refers to a single user query, which can range from scientific inquiries to personal advice requests. Here, the objective is to craft a response $y$ that is both engaging and informative; we make use of the Anthropic Helpful and Harmless dialogue dataset \citep{bai2022training}, which contains 170,000 conversations between humans and automated assistants. Each conversation concludes with two responses from a large-scale (specific model unspecified) language model, each annotated with a human preference indicating the superior reply. As there is no publicly available SFT model tailored for this dataset, we instead fine-tune an existing language model using only those completions that were marked as preferred, resulting in our SFT model for this setting.

\textbf{Evaluation.} We employ two main strategies for evaluating our models. To assess how well each algorithm fulfills the constrained reward maximization objective in the controlled sentiment generation context, we plot the trade-off frontier between the attained reward and the KL-divergence from the reference policy. This analysis is feasible because the true reward function (a sentiment classifier) is accessible. In contrast, when transitioning to real-world conditions—where the actual reward function remains unknown—evaluation shifts to measuring each algorithm’s \textit{win rate} relative to a baseline policy. Here, GPT-4 is utilized as a surrogate for human judgment to compare summary quality in the summarization task and response helpfulness in single-turn dialogue. For the summarization domain, baseline comparisons are conducted using ground-truth reference summaries from the test data; in dialogue, we benchmark against the human-preferred responses contained in the test set. Although prior research suggests that language models often outperform standard metrics as automatic evaluators \citep{Chen2023ExploringTU}, we supplement our analysis with a human evaluation to validate the reliability of GPT-4-based assessments (see Sec.~\ref{sec:human-judgments}). Our findings reveal a high correlation between GPT-4’s judgments and human assessments, with levels of agreement between humans and GPT-4 often matching or exceeding inter-annotator consensus among human evaluators.

\textbf{Methods.} Beyond DPO, we include several prevalent techniques for training language models aligned to human preferences. As a baseline, we apply zero-shot prompting with \textbf{GPT-J} \citep{gpt-j} for the summarization scenario and 2-shot prompting with \textbf{Pythia-2.8B} \citep{biderman2023pythia} for the dialogue experiments. Additionally, we assess models trained with \textbf{SFT} as well as the \textbf{Preferred-FT} variant, which leverages supervised fine-tuning on the preferred output $y_w$, selected from the SFT model in both controlled sentiment and summarization or from a generic language model in the case of single-turn dialogue. Another semi-supervised technique under evaluation is \textbf{Unlikelihood}~\citep{welleck2019neural}, which seeks to increase the likelihood of $y_w$ while simultaneously suppressing the probability of $y_l$, moderated by an adjustable coefficient $\alpha\in[0,1]$ for the 'unlikelihood' penalty. We also investigate \textbf{PPO} \citep{schulman2017proximal}, utilizing a reward model derived from preference data, alongside \textbf{PPO-GT}, an oracle version trained with direct access to the ground truth reward in the controlled sentiment scenario. In the sentiment generation experiments, both an off-the-shelf PPO-GT implementation \cite{leandro_von_werra_2023_7790115} and an enhanced version—with normalized rewards and optimized hyperparameters for improved outcomes—are used; these improvements are similarly applied to PPO when operating with learned rewards. Finally, we consider the \textbf{Best of $N$} approach: $N$ completions are generated from the SFT model (or from Preferred-FT in dialogue), and the one with the top score, according to the learned reward function, is selected. Despite this method’s strong performance—by uncoupling reward model quality from PPO-based optimization—it proves impractical at scale, as inference necessitates sampling $N$ completions for each individual test-time query.

\subsection{How well can DPO optimize the RLHF objective?}

Standard RLHF algorithms optimize a KL-regularized reward maximization objective, which requires balancing the maximization of the reward function with penalization for diverging from the reference policy. Thus, algorithmic comparisons should consider not only the reward attained but also the magnitude of the KL divergence; securing a marginally higher reward at the cost of substantially increased KL is not generally advantageous. Figure~\ref{fig:frontier-tldr-main} depicts the tradeoff frontier between reward and KL across algorithms within the sentiment task. For each algorithm, we carry out several independent runs, varying the policy conservativeness hyperparameter in each instance (target KL $\in\{3,6,9,12\}$ for PPO, $\beta \in \{0.05,0.1,1,5\}$, and $\alpha\in\{0.05,0.1,0.5,1\}$ for unlikelihood, with random initialization for preferred-FT), yielding 22 experiments in total. At intervals of 100 training steps until reaching convergence, policies are evaluated on a common set of held-out prompts by computing both the mean reward, as dictated by the ground-truth reward function, and the mean sequence-level KL\footnote{This is computed as the aggregate sum of KL-divergences across timesteps.} relative to the reference policy $\text{KL}\left(\pi\mid \mid \pi_\text{ref}\right)$. The results reveal that DPO attains a superior reward-KL efficiency, consistently delivering the highest rewards for low KL values. This finding is noteworthy for several reasons: although DPO and PPO are designed to optimize an identical objective, DPO significantly outperforms PPO—its reward/KL frontier strictly subsumes that of PPO. Furthermore, even when PPO has direct access to the true rewards (PPO-GT), DPO continues to surpass PPO’s frontier.

\subsection{Can DPO scale to real preference datasets?}
\label{sec:dpo-real-datasets}
We proceed by assessing the fine-tuning effectiveness of DPO within the contexts of summarization and single-turn dialogue tasks. Traditional automatic metrics like ROUGE often exhibit a weak correlation with human judgment when evaluating summarization quality~\citep{stiennon2022learning}. Earlier research demonstrated that tuning language models with PPO guided by human preferences leads to more favorable summaries. In our experiments, we benchmark multiple approaches by generating responses on the test portion of the TL;DR summarization dataset and determining the mean win rate compared to reference outputs. For each method, samples are drawn across a temperature range from 0.0 to 1.0, and the resulting win rates are depicted in Figure~\ref{fig:frontier-tldr-main} (right). DPO, PPO, and Preferred-FT are each applied to the same GPT-J SFT model\footnote{\url{https://huggingface.co/CarperAI/openai_summarize_tldr_sft}}. Our findings indicate that DPO achieves a win rate of nearly 61\% when using a temperature setting of 0.0, surpassing PPO, which attains about 57\% at its own optimal temperature of 0.0. In addition, DPO secures a superior maximum win rate compared to the best of $N$ strategy. Notably, DPO’s $\beta$ hyperparameter received little tuning, which implies that further optimization could enhance these outcomes. We also observe that DPO maintains robust performance across varying sampling temperatures, in contrast to PPO, whose effectiveness may decline to that of the original GPT-J model at higher temperatures. Preferred-FT, meanwhile, does not yield noticeable improvements over the initial SFT model. Furthermore, we provide a direct comparison between DPO and PPO through human evaluations in Section~\ref{sec:human-judgments}, where DPO samples at a temperature of 0.25 were chosen over PPO’s samples at the same temperature in 58\% of cases.

For single-turn dialogue evaluation, we conduct experiments using the subset of the Anthropic HH dataset's test split \citep{bai2022training} corresponding to a single human-assistant exchange. To assess the different approaches, we rely on GPT-4, which uses the reference preferred completions from the test set to calculate win rates for each method. In the absence of a standard SFT baseline for this scenario, we initialize with a pre-trained Pythia-2.8B model and utilize Preferred-FT to fine-tune a reference model on the preferred completions, ensuring the model-generated responses remain in-distribution, followed by further training with DPO. For comparative purposes, we include the Best of 128 Preferred-FT completions (our analysis revealed that increasing $N$ beyond 128 does not yield further improvements for this dataset; refer to Appendix Figure~\ref{fig:best-of-n}) and a 2-shot prompt-augmented Pythia-2.8B base model. DPO consistently matches or exceeds the top-performing temperatures observed in these baselines. Additionally, we assess an RLHF system—trained with PPO on the Anthropic HH dataset\footnote{\url{https://huggingface.co/reciprocate/ppo_hh_pythia-6B}} and made available by a widely recognized provider\footnote{\url{https://github.com/CarperAI/trlx/tree/main/examples/hh}}—but are unable to identify a prompt or sampling temperature where its performance surpasses that of the unmodified Pythia-2.8B base model. Considering our TL;DR findings and the observation that both approaches optimize for the same reward, we interpret Best of 128 as a reasonable surrogate for PPO-based performance. Ultimately, DPO stands out as the only resource-efficient technique to enhance performance over the dataset's preferred completions, rivaling or surpassing the more computationally intensive Best of 128 baseline. Moreover, as depicted in Figure~\ref{fig:dialogue-main}, DPO reaches optimal performance in relatively few training iterations.

\subsection{Generalization to a new input distribution}

To more thoroughly assess how PPO and DPO respond to distributional changes, we test the policies obtained from our Reddit TL;DR summarization study on a different data source: news articles drawn from the test split of the CNN/DailyMail dataset \citep{nallapati-etal-2016-abstractive}. These evaluations employ the optimal sampling temperatures identified on the TL;DR task (0 and 0.25). Table~\ref{tab:ood} displays the outcomes. For the evaluation metric, we calculate GPT-4’s win rate over ground-truth summaries, using the same GPT-4 (C) prompt from Reddit TL;DR, except that “forum post” is changed to “news article.” On this alternative distribution, DPO continues to substantially outperform PPO. These findings indicate that DPO can generalize at least as robustly as PPO, despite not leveraging the additional unlabeled Reddit TL;DR prompts utilized by PPO.

\subsection{Validating GPT-4 judgments with human judgments}
\label{sec:human-judgments}
A human evaluation study is performed to test the robustness of GPT-4’s judgments, focusing on the TL;DR summarization task and leveraging two distinct GPT-4 prompt variants. The \textbf{GPT-4 (S)} prompt simply asks which summary better encapsulates the essential content of the post. The \textbf{GPT-4 (C)} prompt additionally asks which summary is more concise; this prompt is evaluated because GPT-4 exhibits a tendency to favor more verbose and repetitive outputs when using the \textbf{GPT-4 (S)} prompt, a preference not shared by human annotators. Full prompt texts are provided in Appendix~\ref{app:prompts}. Three head-to-head comparisons are included, featuring the highest-performing (DPO, temp. 0.25), lowest-performing (PPO, temp. 1.0), and an intermediate-performing (SFT, temp. 0.25) method, each compared to the PPO method at its optimal (greedy) temperature to capture a spectrum of summary qualities. Results demonstrate that with either prompt, GPT-4’s selections are about as consistent with human assessments as human raters are with each other, indicating that GPT-4 serves as a suitable proxy for human preferences (multiple human ratings were only collected for DPO and PPO-1 comparisons due to the limited rater pool). In aggregate, the \textbf{GPT-4 (C)} prompt yields win rates that most closely match human evaluations, and thus it is adopted for the main experiments in Section~\ref{sec:dpo-real-datasets}. Further methodological specifics—including the web-based rating interface and details on participating raters—are described in Appendix~\ref{app:human-study}.

\section{Discussion}
Preference-based learning offers a robust and scalable approach for developing language models that are both capable and aligned with human intentions. In this work, we have presented DPO—a straightforward training methodology that enables language models to learn from preferences without relying on reinforcement learning frameworks. Unlike methods that reshape preference learning to fit the mold of conventional RL in order to utilize typical RL algorithms, DPO establishes a direct correspondence between language model policies and reward functions. This mapping permits the training of language models to conform to human preferences through a straightforward cross-entropy loss function, circumventing the need for reinforcement learning without sacrificing generality. Remarkably, with minimal hyperparameter adjustment, DPO achieves results on par with, or superior to, leading RLHF methods such as those employing PPO. As a result, DPO significantly lowers the barriers associated with training language models using human preferences.

\textbf{Limitations \& Future Work.} Our findings open several avenues for continued investigation. One key question is how the DPO-trained policy performs on out-of-distribution data relative to approaches that learn from an explicit reward function. Preliminary findings indicate that DPO models may generalize comparably to those trained with PPO, but a more thorough assessment is warranted. For instance, could self-labeling with DPO be leveraged to take advantage of unlabeled prompts effectively? Another open question concerns the manifestation of reward over-optimization in the direct preference optimization framework; the moderate dip in performance observed in Figure~\ref{fig:dialogue-main}-right might exemplify this phenomenon. Moreover, our evaluations thus far have been limited to models containing up to 6B parameters; scaling DPO to substantially larger, state-of-the-art models represents a promising research direction. Regarding assessment strategies, we note that GPT-4's computed win rates are influenced by prompt phrasing; subsequent research might explore optimal techniques for eliciting reliable evaluations from automated systems. Lastly, the scope of DPO extends well beyond language modeling from human preferences—it has potential applications for training generative models across various modalities.